\question{}

While inspecting the cluster, you discover that one of your colleagues Alice has
been mining bitcoin on the compute nodes whenever no user job is running on them,
i.e.\ he keeps the otherwise-idle nodes busy at $100\%$ utilization.

Bob complains that this is inflating the electricity bill.
Alice replies:
\begin{quote}\itshape
``The nodes have to stay powered on anyway as they must be ready the instant a user
asks for compute, so we are not allowed to switch them off. A node that is on
consumes its power regardless of what it does, so it draws the same whether it sits
idle or mines bitcoin. I might as well put those idle cycles to use.''
\end{quote}

We grant Alice her premise: the nodes \textbf{must stay on and ready}, so the question is
\emph{not} whether they could be switched off.
Taking that as fixed, is the rest of
Alice reasoning ``that an \textbf{on-but-idle} node draws the same power as an
\textbf{on-and-fully-loaded} one'' correct? Justify quantitatively.

\SetTotalwidth
  \begin{solutionbox}{14cm}
    \textbf{No --- her reasoning is wrong; ``on'' does not mean ``drawing full power''.}
    The power of a CPU/GPU is not a constant; it splits into
    \[
      P \;=\; \underbrace{P_{\text{static}}}_{\text{leakage}} \;+\; \underbrace{\alpha\,C\,V^2 f}_{\text{dynamic}} ,
    \]
    where $\alpha$ is the \emph{activity factor} (fraction of transistors actually
    switching). On an idle node $\alpha \approx 0$, so the dynamic term --- which
    dominates the total under load --- essentially vanishes. On top of that, an idle
    node is not held at its maximum operating point:
    \begin{itemize}
      \item \textbf{DVFS} lowers the clock $f$ and voltage $V$ when there is no work;
            since dynamic power $\propto V^2 f$, dropping both cuts it
            super-linearly.
      \item \textbf{Idle / sleep states} (CPU C-states, GPU power states) clock-gate
            and power-gate unused units, further shrinking $P_{\text{static}}$.
    \end{itemize}
    Concretely, a server CPU/GPU idles at a small fraction of its TDP and rises to
    (near) its full TDP only at $100\%$ utilization --- often a $3\times$ to $10\times$
    swing between idle and a sustained mining load. Bitcoin mining deliberately pins
    every core at maximum, i.e.\ it drives $\alpha$, $f$ and $V$ all to their peak, so
    it adds (close to) the \emph{entire} dynamic power budget of every node it touches
    --- plus extra cooling power to remove that heat.

    So even though the nodes legitimately stay powered on, ``idle'' and ``mining'' are
    very far from the same electricity draw: the mining is genuinely adding to the bill,
    and the ``it's the same anyway'' argument is false.
  \end{solutionbox}
